\section{Method}
\label{sec:methodology}

This section formalizes the contract-coverage optimization problem as a
Markov decision process (MDP) and describes the reinforcement-learning
algorithm we adopt, Proximal Policy Optimization (PPO). We restrict the
formulation to the exact scope of our empirical study: a \emph{single
piece} of medical equipment followed over a finite daily horizon. This
deliberate simplification makes the analysis tractable and keeps the
theoretical description and the implementation in strict one-to-one
correspondence; extensions to a portfolio of heterogeneous equipment are
discussed in Section~\ref{sec:discussion}.


\subsection{MDP formulation}
\label{subsec:mdp}

We model the problem as a finite-horizon MDP
$\mathcal{M} = (\mathcal{S}, \mathcal{A}, P, r, T)$, with daily time steps
and horizon $T = 120$ days, tracking a single piece of equipment.


\paragraph{State space.}
The state at step $t$ is a four-dimensional vector
\[
    s_t = (\mathrm{covered}_t,\;
           \mathrm{days\_to\_exp}_t,\;
           \mathrm{cost\_level},\;
           \mathrm{risk\_level}),
\]
where $\mathrm{covered}_t \in \{0, 1\}$ indicates whether an active
contract covers the equipment, $\mathrm{days\_to\_exp}_t$ is the remaining
number of days before expiration of the current coverage,
$\mathrm{cost\_level} \in \{0, 1, 2\}$ discretizes the contractual cost
(low, medium, high), and $\mathrm{risk\_level} \in \{0, 1, 2\}$
discretizes the equipment risk class (low, medium, high). The two
discrete attributes are sampled at the start of each episode and remain
constant throughout. All four components are normalized to the unit
interval before being fed to the policy.


\paragraph{Action space.}
The action space is discrete with three elements:
\[
    \mathcal{A} = \{\, 0: \text{do nothing},\;
                      1: \text{renew contract},\;
                      2: \text{extend warranty} \,\}.
\]
A \emph{renewal} (action 1) resets the remaining coverage to a fixed
duration ($\mathrm{RENEWAL\_DAYS} = 30$) and restores the covered status
if the equipment was uncovered. A \emph{warranty extension} (action 2)
adds a fixed duration ($\mathrm{EXTENSION\_DAYS} = 30$, capped at
$\mathrm{MAX\_COVERED\_DAYS} = 60$) at a reduced cost, but requires an
already-active contract; invoking it while uncovered incurs a small
no-op penalty.


\paragraph{Reward function.}
The per-step reward aggregates six components, each chosen so that the
total remains in a numerically stable range and induces a clear ordering
between desirable and undesirable decisions:
\begin{align*}
r_t \;=\; &\; r^{\text{coverage}}_t
          + r^{\text{timing}}_t
          + r^{\text{cost}}_t
          + r^{\text{expiry}}_t
          + r^{\text{early}}_t
          + r^{\text{failure}}_t.
\end{align*}
The continuous coverage term is $+\mathrm{COVERAGE\_BONUS}$ while covered
and $-\mathrm{UNCOVERED\_PENALTY}$ otherwise. The timing term rewards
renewals or extensions issued close to expiration
($+\mathrm{RENEWAL\_CLOSE\_BONUS}$ when
$\mathrm{days\_to\_exp}_t \le 5$) and penalizes early renewals
($-\mathrm{RENEWAL\_EARLY\_PENALTY}$ when
$\mathrm{days\_to\_exp}_t \ge 20$). The cost term is a contract cost
proportional to $\mathrm{cost\_level}$ (smaller for extensions than for
full renewals). The expiration term is a one-shot penalty applied the
exact step a contract expires without renewal. The failure term is a
one-shot penalty whenever a failure occurs while the equipment is
uncovered, sampled with probability proportional to
$(1 + \mathrm{risk\_level})$. Numerical coefficients are listed in
Section~\ref{subsec:training}.


\paragraph{Transition dynamics.}
At each step, the environment sequentially (i) applies the chosen action,
(ii) decrements $\mathrm{days\_to\_exp}_t$ by one and triggers expiration
if the counter reaches zero, (iii) samples a Bernoulli failure indicator
with risk-conditional probability, and (iv) emits the scalar reward. The
static attributes $\mathrm{cost\_level}$ and $\mathrm{risk\_level}$ are
unchanged across a full episode, so the only stochastic transitions
concern the coverage counter and the failure process. The dynamics are
fully Markovian in the four-dimensional state.


\subsection{Tabular Bellman optimum (value iteration)}
\label{subsec:vi-baseline}

For benchmarking and clarification of optimality notions, we also compute an
exact greedy policy in a finite MDP that augments each physical state with the
elapsed time index $\tau \in \{0,\dots,T{-}1\}$ recorded by the simulator at
each decision (the environment's internal step counter before the action is
applied).

Because $\mathrm{cost\_level}$ and $\mathrm{risk\_level}$ are fixed throughout
an episode, this augmented process is finite---at most
$|\mathcal{S}_\text{tab}| \approx 65{,}880$ reachable tuples for our
$T{=}120$, $D_{\max}{=}60$ setting---so we solve it with synchronous
discounted value iteration using the same $\gamma{=}0{.}99$ as PPO.

The uncovered failure channel affects only stochastic roll-outs; Bellman backups
evaluate its contribution through the conditional expected penalty given the
coverage status after deterministic transition of the expiry counter. The
associated greedy policy, denoted \emph{VI (exp.)} in experiments, maximises the
usual discounted criterion in \emph{expectation} relative to this tabulation;
Monte Carlo evaluation on the sampled simulator inherits Bernoulli noise, so confidence
intervals remain non-degenerate.


\subsection{Proximal Policy Optimization (PPO)}
\label{subsec:ppo}

To approximate the MDP defined in Section~\ref{subsec:mdp}, we use the
\emph{Proximal Policy Optimization} (PPO)
algorithm~\cite{schulman2017ppo}, a policy-gradient method recognized for
its empirical stability and implementation simplicity.


\paragraph{Principle.}
PPO belongs to the family of actor-critic methods. The algorithm jointly
trains two components: a stochastic policy $\pi_\theta(a \mid s)$ that
maps each state to a distribution over the three admissible actions, and
a value function $V_\phi(s)$ that estimates the expected future return
from state $s$. Both components are parameterized by neural networks
sharing a common two-layer encoder (64 units, $\tanh$ activations) with
separate softmax policy and scalar value heads. They are learned from
trajectories collected under the current policy. The advantage estimate,
which measures the marginal gain of an action relative to the average
value of the state, is obtained through Generalized Advantage Estimation
(GAE)~\cite{schulman2016gae}, with a discount $\gamma = 0.99$ and GAE
coefficient $\lambda = 0.95$, providing a controllable trade-off between
bias and variance.


\paragraph{Clipped objective.}
The specificity of PPO lies in its learning objective. Let
$r_t(\theta) = \pi_\theta(a_t \mid s_t) / \pi_{\theta_{\text{old}}}(a_t \mid s_t)$
denote the likelihood ratio between the current policy and the behavior
policy that produced the trajectory, and let $\hat{A}_t$ denote the
normalized GAE advantage. PPO maximizes
\[
    \mathcal{L}^{\text{clip}}(\theta)
    = \mathbb{E}_t \Big[
        \min \big(r_t(\theta)\, \hat{A}_t,\;
                 \mathrm{clip}(r_t(\theta),\, 1 - \epsilon,\, 1 + \epsilon)\, \hat{A}_t
            \big)
      \Big],
\]
with clip parameter $\epsilon = 0.2$. The value function is trained by
minimizing the squared return-prediction error, and an entropy bonus is
added to the global objective to encourage exploration and to prevent
premature convergence to a deterministic policy.


\paragraph{Implementation details.}
Observations are normalized online through a running estimate of their
mean and variance (Welford's algorithm) and clipped to $[-5, 5]$ before
being fed to the network. Advantages are normalized per batch. At each
iteration, a full episode is collected under the current policy, and the
policy--value pair is updated for four epochs over mini-batches of size
$64$. A complete list of hyperparameters is given in
Section~\ref{subsec:training}. Training is conducted entirely offline on
the simulator, which enables intensive exploration without operational
risk and yields a reproducible evaluation of the learned policy prior to
any deployment.
